\chapter{Coordination, Its Absence, and Partial Remedies}
\label{ch:coordination}

\begin{quotation}
\noindent\textit{Synopsis.} This chapter investigates what happens when
stages within a relay attempt to compensate for the absence of global
coordination. Augmentation, provenance, priors, adjudication, and revision
are reinterpreted not as an unrelated list of techniques but as distinct
strategies for partially coordinating otherwise independent repair
economies, each addressing a different failure mode: augmentation where
hidden state is enumerable, provenance where loss can at least be recorded,
priors where a consumer can plausibly fill a benign gap, adjudication where
authority can terminate ambiguity without recovering information, and
revision where the relay itself must be replaced. None fully solves the
coordination problem; each narrows a specific failure mode.
\end{quotation}

\section{Why Full Coordination Is Usually Unavailable}

A field-recording engineer, the vendor of a transcription model, and a
human editor who later summarizes the transcript typically share no
protocol, no common budget, and often no direct communication at all.
\cref{def:budget-relay}'s uncoordinated case is not a pathological edge
case to be assumed away; for most substrate chains of practical interest,
it is the default condition. This chapter asks what a stage \emph{can} do
to compensate for that default, given that Chapter~\ref{ch:synchronic-to-diachronic}
already showed (\cref{prop:no-global-conserved-quantity}) that no
extension of Chapter~\ref{ch:repair-economics}'s single-observer
apparatus is available to fall back on.

\begin{definition}[Coordination Mechanism]
\label{def:coordination-mechanism}
A coordination mechanism is any deliberate modification to
\cref{def:budget-relay}'s default independence condition, permitting some
stage or external party access to information --- a loss ledger, a
hidden-state enumeration, an authoritative ruling, or the raw root object
itself --- that a fully uncoordinated relay would not provide.
\end{definition}

Five such mechanisms recur across the domains treated in
Part~\ref{part:worked-domains}, and it is worth resisting the temptation
to treat them as an unrelated list of ad hoc devices collected because
they happen to be useful. Each addresses a specific, identifiable gap in
\cref{def:budget-relay}'s independence condition, and organizing them by
which gap they close is more informative than organizing them
alphabetically or by domain.

\section{Augmentation}
\label{sec:coord-augmentation}

Augmentation adds an enumerated hidden-state set $H$ directly to a
rendering, pre-committing budget to distinctions a stage's own task does
not itself require, on the strength of knowing in advance that some later
consumer will. By \cref{prop:sufficiency-closed-only}, this strategy is
viable only in closed hidden-state domains (\cref{def:closed-open}):
Chapter~\ref{ch:chess} develops the case in full, where $H$ --- castling
rights, en passant status, repetition count --- is finite and known at
the time a position is recorded, so that funding it requires no
prediction about which future task will need it, only the (comparatively
modest) foresight that \emph{some} task among a known, finite list might.

\begin{remark}[Augmentation Does Not Solve the Open Case]
\label{rem:augmentation-limit}
Augmentation is unavailable, not merely difficult, in open hidden-state
domains (\cref{def:closed-open}): there is no finite $H$ to pre-commit
budget to, since \cref{prop:sufficiency-closed-only} shows no finite
augmentation closes an open domain's hidden-state space. Painting
(Chapter~\ref{ch:painting}) is not a case where augmentation would help
if only more budget were available; it is a case where augmentation, as a
strategy, does not apply regardless of budget.
\end{remark}

\section{Provenance and Loss-Ledger Sharing}
\label{sec:coord-provenance}

\begin{definition}[Provenance-Coordinated Relay]
\label{def:provenance-coordinated}
A substrate chain is provenance-coordinated if \cref{def:budget-relay}'s
independence condition is relaxed to permit stage $i$'s transition $B_i$
to be selected with reference to the loss ledgers $L_0, \ldots, L_{i-1}$
of every earlier stage, while still forbidding reference to any later
stage's budget or ledger.
\end{definition}

This is the weakest relaxation of \cref{def:budget-relay} that still
qualifies as coordination at all: it requires no advance agreement
between stages, only that each ledger be produced and made accessible to
whatever comes after it, and it preserves the chain's directionality ---
a stage still cannot see the future, only more of the past than
\cref{def:budget-relay} otherwise allows.

\begin{proposition}[Provenance Sharing Does Not Restore Sufficiency Alone]
\label{prop:provenance-insufficient}
A provenance-coordinated relay (\cref{def:provenance-coordinated}) can
still exhibit the failure mode of \cref{thm:composition-failure}.
\end{proposition}

\begin{proofsketch}
Reading $L_0, \ldots, L_{i-1}$ tells stage $i$, and eventually a
downstream consumer, \emph{what} was discarded and under what budget ---
but not what a future, unknown $\tau^*$ will require, since
\cref{prop:task-unavailability} applies to the ledger's reader exactly as
it applies to the stage that produced the ledger. Knowing that $d^*$ was
dropped at stage $0$ does not, by itself, supply $d^*$; recovering it
requires either \cref{prop:root-recovery}'s root access, which
\cref{def:provenance-coordinated} does not grant, or advance knowledge of
$\tau^*$, which \cref{prop:task-unavailability} rules out in general.
Provenance sharing converts an undetectable insufficiency into a
diagnosable one; it does not convert it into an avoided one.
\end{proofsketch}

\cref{prop:provenance-insufficient} is why Chapter~\ref{ch:budget-relay}
described the budget-aware loss ledger (\cref{def:loss-ledger}) as a
coordination device rather than a solution: it is necessary for a
downstream party to \emph{diagnose} what happened, which matters a great
deal for the reasons Chapter~\ref{ch:pipelines} develops, but it is not
sufficient to \emph{prevent} what happened, which requires the stronger
and costlier machinery of root preservation.

\section{Priors}
\label{sec:coord-priors}

A consumer-side substitute for coordination: when a relay supplies no
interference information, the consumer supplies plausible reconstructions
of the missing distinction from general knowledge. Painting
(\cref{ch:painting}) is the paradigm case, and it is available precisely
because, and only because, the discarded distinction's absence is benign
(\cref{def:benign-loss}) relative to the consumer's task: the correctness
set $\mathcal{C}_\tau(e)$ of \cref{def:correctness-set} is large enough
that a plausible guess, rather than the exact discarded distinction, is
sufficient to land inside it. Priors are unavailable, or worse than
unavailable, exactly where \cref{def:corrupting-loss}'s conditions hold:
a confidently wrong prior substituted for an undisclosed loss is
corruption compounded by false confidence, not a benign completion.

\section{Adjudication}
\label{sec:coord-adjudication}

\begin{proposition}[Adjudication Terminates Ambiguity]
\label{prop:adjudication}
Let $\mathcal{J}$ be an authority external to a substrate chain.
Adjudication by $\mathcal{J}$ produces a unique continuation even when the
distinguishability information available in the chain is incomplete.
\end{proposition}

\begin{proofsketch}
Under adjudication, the continuation is selected by an authoritative
decision procedure rather than inferred from available evidence. The
underlying information deficit --- the discarded distinctions --- is
unaffected by the ruling; what changes is that ambiguity is terminated
procedurally rather than resolved informationally.
\end{proofsketch}

It is worth stating precisely, in the vocabulary of
Chapter~\ref{ch:admissibility}, what adjudication actually does to
\cref{def:correctness-set}'s correctness set. Adjudication does not
discover which continuation truly belongs to $\mathcal{C}_\tau(e)$; it
substitutes $\mathcal{J}$'s ruling as the operative correctness set going
forward, regardless of whether that ruling coincides with whatever the
undiscoverable true $\mathcal{C}_\tau(e)$ would have contained.
Chapter~\ref{ch:contracts} treats this substitution as the defining
feature of legal continuation under irreducible uncertainty: courts do
not typically claim to have recovered the parties' true intent so much as
to have settled, bindingly, what will be treated as though it were.
Adjudication is, in this taxonomy, the only mechanism that terminates
ambiguity without requiring any additional distinguishability information
to be recovered or supplied. It trades informational completeness for a
binding decision.

\section{Revision}
\label{sec:coord-revision}

\begin{definition}[Relay Revision]
\label{def:relay-revision}
Wholesale replacement of an entire substrate chain $A_0 \to \cdots \to
A_n$ by a new chain built from new production, rather than re-derived
from the original chain's root object.
\end{definition}

Revision is easy to confuse with root recovery (\cref{prop:root-recovery}),
and it is worth being precise about the difference. Root recovery
re-derives a new downstream rendering $A_i'$ from the \emph{same}
preserved $A_0$, under a revised discard policy; it presupposes $A_0$
survived and is available. Revision, as defined here, does not presuppose
this: it discards the entire existing chain, root included, and begins
again from new evidence, a new theory, or a new production process
entirely. Chapter~\ref{ch:science} develops this case in full: a
superseded scientific theory is not re-derived from some preserved
``root'' observation set in the way \cref{prop:root-recovery} re-derives
a rendering; the theory itself, and often the observational framework
that produced the evidence for it, is replaced.

\section{Coordination Versus Pooling: A Boundary Case}
\label{sec:coord-vs-pooling}

The five mechanisms above all operate \emph{within} the assumption that
stages remain separate --- each retains its own local budget, and
coordination consists of granting limited, structured access across that
separation (a ledger here, an authority there), not erasing it.
Appendix~\ref{app:comparison}'s demonstration that pooling the witness's
stage budgets into one shared total can restore sufficiency
(\cref{prop:pooling-restores}) is a more radical move than any mechanism
in this chapter's taxonomy: pooling does not coordinate separate local
economies, it eliminates the separation itself, collapsing the relay back
into the single-observer setting Chapter~\ref{ch:repair-economics}
already covers. None of augmentation, provenance-sharing, priors,
adjudication, or revision does this; each leaves
\cref{def:substrate-chain}'s stage structure intact and works within it.
This is worth flagging explicitly because pooling is, in a precise sense,
strictly more powerful than anything on this chapter's list --- it is
what makes \cref{thm:redistribution} available again --- and its absence
from the taxonomy is not an oversight but a reflection of how rarely it
is actually achievable: pooling requires exactly the kind of unified,
advance-negotiated authority over an entire chain's resources that
Section~\ref{sec:coord-augmentation}'s discussion of augmentation already
identified as the exception rather than the rule.

\section{The Limits of the Taxonomy}

\begin{proposition}[No Coordination Mechanism Restores Universal Sufficiency]
\label{prop:no-coordination-universal}
No coordination mechanism drawn from
Sections~\ref{sec:coord-augmentation}--\ref{sec:coord-revision} can
guarantee sufficiency of a substrate chain over an open task space $T$.
\end{proposition}

\begin{proofsketch}
This follows directly from \cref{prop:task-unavailability}: each
coordination mechanism above operates by preserving, recovering, or
authoritatively resolving \emph{some additional} set of distinctions, but
$T$'s openness means a task $\tau^*$ can always be constructed requiring a
distinction outside whatever finite additional set any given mechanism
supplies. Coordination narrows the gap identified in
\cref{principle:global-insufficiency-restated}; it does not close it. This
result was labeled a Theorem in an earlier draft; the label has been
corrected to Proposition here, consistent with the monograph's
epistemic-labeling discipline, since the argument is a direct application
of \cref{prop:task-unavailability} rather than an independent derivation.
\end{proofsketch}

Chapter~\ref{ch:composition-failure} is, in effect, the sharper and more
constructive version of \cref{prop:no-coordination-universal}: it does
not merely state that no mechanism guarantees universal sufficiency, but
exhibits the mechanism by which locally rational, uncoordinated stages
compose into global insufficiency in the first place, and
Appendix~\ref{app:proof} exhibits it with a construction robust to the
choice of local discard rule.
